% !TeX program = pdflatex
\documentclass[11pt]{extarticle}

\usepackage[T2A]{fontenc}
\usepackage[utf8]{inputenc}
\usepackage[english,russian]{babel}
\usepackage{amsmath,amssymb,amsfonts}
\usepackage{mathtools}
\usepackage{hyperref}
\usepackage{geometry}
\geometry{a4paper, left=1.5cm, right=1.5cm, top=2cm, bottom=2cm}
\usepackage{setspace}
\onehalfspacing
\usepackage{indentfirst}
\usepackage{cite}
\usepackage{xcolor}

\begin{document}

\begin{center}
    \Large\textbf{Улучшение оценки нормы разности гессианов} \\
    \vspace{0.5cm}
    \large Игнатьев Даниил и Никита Киселёв  \\
    % \normalsize Институт проблем передачи информации им. А.А. Харкевича РАН, Москва \\
    \texttt{ignatev.daniil@phystech.edu}
\end{center}

\vspace{1cm}

% Аннотация: заголовок по центру, текст по центру в minipage
\begin{center}
    \textbf{Аннотация}
\end{center}
\begin{center}
    \small
    %\tiny
    \begin{minipage}{0.85\textwidth}
        Исследование локальных свойств поверхности функции потерь, в частности поведения матрицы Гессе, является ключевой задачей для понимания сходимости нейросетей и оценки необходимого объема данных. В существующих работах оценка сходимости сводится к анализу нормы разности между матрицей Гессе на новом объекте и средней матрицей Гессе по предыдущей выборке. Текущий подход использует для этой оценки простое неравенство треугольника, что дает грубую константную верхнюю границу $2M_H$, не зависящую от размера выборки и не отражающую интуитивно ожидаемого убывания ошибки при росте $k$. Таким образом, проблема заключается в отсутствии теоретической оценки, которая описывала бы скорость убывания этой нормы.

        Целью данной работы является вывод более точной верхней границы для нормы разности матриц Гессе, которая будет учитывать эффект усреднения и убывать с ростом количества наблюдений $k$. Для этого используется разложение разности через популяционное среднее и вводятся предположения о концентрации пер-объектных матриц Гессе. Основным результатом станет получение оценки вида $O(1/k)$ или $O(1/\sqrt{k})$ в зависимости от предположений с явными константами. Полученная теоретическая оценка будет подтверждена в экспериментах на задачах классификации изображений (MNIST, FashionMNIST, CIFAR10). Значимость исследования заключается в том, что улучшенная оценка позволит дать более строгое обоснование скорости сходимости оптимизационной поверхности, что является важным шагом для развития методов анализа достаточного размера обучающей выборки в глубоком обучении.

        \textbf{Ключевые слова:} матрица Гессе, спектральная норма, неравенства концентрации, нейронные сети, размер выборки.
    \end{minipage}
\end{center}

\vspace{1cm}

\section{Введение}

В статистической теории машинного обучения \cite{Vapnik1995} качество модели напрямую зависит от размера обучающей выборки. Для параметрических моделей функция потерь вычисляется по ограниченному набору данных, поэтому возникает вопрос о том, как определить оптимальный размер обучающей выборки. Классические подходы к оценке достаточного объема данных опираются на вероятностные предположения о распределении данных и параметрах модели \cite{Motrenko2014, Bishop2006}, что часто труднопроверяемо на практике. В то же время эмпирические законы масштабирования \cite{Kaplan2020, Hoffmann2022} демонстрируют устойчивые степенные зависимости между объемом данных, числом параметров и качеством модели, однако их теоретическое обоснование остается открытой проблемой.

В последние годы активно исследуется связь между размером выборки и геометрией оптимизационной поверхности — ландшафта функции потерь в пространстве параметров \cite{Li2018, Fort2019}. Локальные свойства ландшафта, такие как кривизна в окрестности минимума, определяются матрицей Гессе. Спектральные характеристики этой матрицы позволяют судить о стабилизации поверхности при увеличении объема данных \cite{Papyan2019, Ghorbani2019, Granziol2021}. В работах \cite{Kiselev2024, Meshkov2024} предложены критерии сходимости ландшафта, основанные на анализе разности эмпирических функций потерь на выборках разного размера. В частности, для оценки изменения ландшафта используется норма разности матриц Гессе на новом объекте и среднего по предыдущим объектам. Текущая верхняя граница этой нормы, полученная с помощью неравенства треугольника, имеет вид \(2M_H\) и не отражает ожидаемого убывания при росте числа наблюдений \(k\). Таким образом, возникает задача получения более точной оценки, учитывающей эффект усреднения и, возможно, вероятностную природу гессианов.

Целью настоящей работы является улучшение теоретической верхней границы для нормы разности матриц Гессе при увеличении размера выборки, отражающей убывание с ростом \(k\). Для достижения этой цели будут использованы вероятностные неравенства концентрации для случайных матриц, а также результаты о спектральных свойствах гессианов в нейронных сетях \cite{Ghorbani2019, Granziol2021}. Полученная оценка позволит уточнить скорость сходимости ландшафта и приблизить решение задачи определения достаточного объема данных.

\section{Постановка задачи}

Рассмотрим задачу обучения параметрической модели \(f_{\mathbf{w}}: \mathbb{X} \to \mathbb{Y}\) с функцией потерь \(\ell(\hat{y}, y)\). Пусть \(\mathfrak{D}_k = \{(\mathbf{x}_i, y_i)\}_{i=1}^k\) — обучающая выборка объема \(k\), составленная из независимых наблюдений. Эмпирическая функция потерь определяется как
\[
\mathcal{L}_k(\mathbf{w}) = \frac{1}{k}\sum_{i=1}^k \ell(f_{\mathbf{w}}(\mathbf{x}_i), y_i).
\]
Обозначим через \(\mathbf{H}_i(\mathbf{w}) = \nabla^2_{\mathbf{w}} \ell(f_{\mathbf{w}}(\mathbf{x}_i), y_i)\) матрицу Гессе на \(i\)-м объекте, а через \(\mathbf{H}^{(k)}(\mathbf{w}) = \frac{1}{k}\sum_{i=1}^k \mathbf{H}_i(\mathbf{w})\) — усредненную матрицу Гессе на выборке.

В работах \cite{Kiselev2024, Meshkov2024} при анализе сходимости ландшафта функции потерь возникает величина
\[
\mathbf{D}_k = \mathbf{H}_{k+1} - \mathbf{H}^{(k)},
\]
которая характеризует изменение кривизны при добавлении нового объекта. Существующая оценка использует неравенство треугольника и предположение об ограниченности спектральной нормы гессианов: если \(\|\mathbf{H}_i\|_2 \le M\) для всех \(i\), то
\[
\|\mathbf{D}_k\|_2 \le \|\mathbf{H}_{k+1}\|_2 + \|\mathbf{H}^{(k)}\|_2 \le 2M.
\]
Эта оценка является детерминированной и не зависит от \(k\), однако она часто оказывается завышенной, поскольку не учитывает вероятностную природу гессианов и эффект усреднения.

Целью данной работы является получение вероятностной верхней границы для \(\|\mathbf{D}_k\|_2\), которая с заданной вероятностью \(1-\delta\) даёт значение, меньшее \(2M\), и, возможно, убывает с ростом \(k\) за счёт члена, связанного с отклонением среднего.

\subsection{Предположения}

Для формулировки результатов введём следующие предположения:

\begin{enumerate}
    \item \textbf{Независимость и одинаковая распределённость:} матрицы \(\mathbf{H}_1, \mathbf{H}_2, \dots\) независимы и одинаково распределены как случайные элементы в пространстве симметричных матриц размера \(N \times N\).
    \item \textbf{Ограниченность:} существует константа \(M > 0\) такая, что \(\|\mathbf{H}_i\|_2 \le M\) почти наверное. (Это стандартное предположение, используемое в диссертации.)
    \item \textbf{Существование среднего:} определим популяционное среднее \(\bar{\mathbf{H}} = \mathbb{E}[\mathbf{H}_i]\). Заметим, что \(\|\bar{\mathbf{H}}\|_2 \le M\).
    \item \textbf{Дополнительные моментные условия:} для применения концентрационных неравенств достаточно ограниченности, но можно также использовать информацию о дисперсии (например, \(\mathbb{E}\|\mathbf{H}_i - \bar{\mathbf{H}}\|_2^2 \le \sigma^2\)).
\end{enumerate}

% Далее мы рассмотрим два случая, различающихся выбором точки, в которой вычисляются гессианы.

% \vspace{0.5cm}
% % Синий вариант: фиксированная точка w_*
% \textcolor{blue}{%
% \subsection*{Случай 1: Фиксированная точка \(\mathbf{w}_*\)}
% Пусть \(\mathbf{w}_*\) — некоторая фиксированная точка в пространстве параметров, не зависящая от обучающей выборки (например, минимум популяционного риска или любая другая точка, выбранная априори). Везде далее вычисления производятся в этой точке, поэтому аргумент \(\mathbf{w}_*\) будем опускать: \(\mathbf{H}_i = \mathbf{H}_i(\mathbf{w}_*)\), \(\mathbf{H}^{(k)} = \mathbf{H}^{(k)}(\mathbf{w}_*)\)

% В этом случае матрицы \(\mathbf{H}_i\) независимы и одинаково распределены, и для разности
% \[
% \mathbf{D}_k = \mathbf{H}_{k+1} - \mathbf{H}^{(k)}
% \]
% справедливо представление
% \[
% \mathbf{D}_k = (\mathbf{H}_{k+1} - \bar{\mathbf{H}}) - \frac{1}{k}\sum_{i=1}^k (\mathbf{H}_i - \bar{\mathbf{H}}).
% \]
% Применяя матричные концентрационные неравенства (например, неравенство Хёфдинга для случайных матриц \cite{Hoeffding1963} или неравенство Бернштейна \cite{Bernstein1946}), можно показать, что с вероятностью не менее \(1-\delta\) выполняется
% \[
% \|\mathbf{D}_k\|_2 \le C_1\sqrt{\frac{\ln(1/\delta)}{k}} + C_2\frac{\ln(1/\delta)}{k},
% \]
% где константы \(C_1, C_2\) зависят от \(M\) и, возможно, от дисперсии \(\sigma^2\). Таким образом, в этом случае оценка убывает с ростом \(k\) как \(O(1/\sqrt{k})\) (или \(O(1/k)\) при дополнительных условиях).
% }

% \vspace{0.5cm}
% % Красный вариант: точка минимума эмпирического риска w_k^*
% \textcolor{red}{%
% \subsection*{Случай 2: Точка минимума эмпирического риска \(\mathbf{w}_k^*\)}
% Пусть \(\mathbf{w}_k^*\) — точка минимума эмпирического риска \(\mathcal{L}_k(\mathbf{w})\). В этом случае гессианы \(\mathbf{H}_i(\mathbf{w}_k^*)\) зависят от всей выборки через \(\mathbf{w}_k^*\), что нарушает независимость и одинаковую распределённость. Тем не менее, интерес представляет величина
% \[
% \mathbf{D}_k = \mathbf{H}_{k+1}(\mathbf{w}_k^*) - \mathbf{H}^{(k)}(\mathbf{w}_k^*),
% \]
% которая возникает при анализе сходимости ландшафта в работе \cite{Kiselev2024}.

% Для получения вероятностной оценки в этом случае требуется учитывать дополнительную ошибку, связанную с отклонением \(\mathbf{w}_k^*\) от истинного минимума. Используя результаты о smoothness функции потерь и концентрации гессианов в окрестности минимума, можно ожидать оценку вида
% \[
% \|\mathbf{D}_k\|_2 \le C_1\sqrt{\frac{\ln(1/\delta)}{k}} + C_2\frac{\ln(1/\delta)}{k} + C_3\|\mathbf{w}_k^* - \mathbf{w}_*\|_2,
% \]
% где последний член учитывает смещение из-за замены точки. При выполнении условий регулярности и достаточно быстрой сходимости \(\mathbf{w}_k^*\) к \(\mathbf{w}_*\) (например, со скоростью \(O(1/\sqrt{k})\)) общий порядок может сохраниться как \(O(1/\sqrt{k})\), но с большей константой.
% }

% \vspace{0.5cm}

% В обоих случаях итоговая цель — получить явные выражения для констант и подтвердить их экспериментально на реальных данных.

\begin{thebibliography}{99}

\bibitem{Vapnik1995} Vapnik V. N. The Nature of Statistical Learning Theory. -- Springer, 1995.
\bibitem{Kaplan2020} Kaplan J. et al. Scaling laws for neural language models // arXiv:2001.08361. 2020.
\bibitem{Hoffmann2022} Hoffmann J. et al. Training compute-optimal large language models // NeurIPS. 2022.
\bibitem{Li2018} Li H. et al. Visualizing the loss landscape of neural nets // NeurIPS. 2018.
\bibitem{Fort2019} Fort S., Ganguli S. Emergent properties of the local geometry of neural loss landscapes // arXiv:1910.05929. 2019.
\bibitem{Papyan2019} Papyan V. Measurements of three-level hierarchical structure in the outliers in the spectrum of deepnet Hessians // ICML. 2019.
\bibitem{Motrenko2014} Motrenko A., Strijov V. Sample size estimation for logistic regression using cross-validation and Kullback–Leibler divergence // JMLR. 2014.
\bibitem{Bishop2006} Bishop C. M. Pattern Recognition and Machine Learning. -- Springer, 2006.
\bibitem{Ghorbani2019} Ghorbani B., Krishnan S., Xiao Y. An investigation into neural net optimization via Hessian eigenvalue density // ICML. 2019.
\bibitem{Granziol2021} Granziol D., Zohren S., Roberts S. Hessian Eigenspectra of More Realistic Nonlinear Models // NeurIPS. 2021.
\bibitem{Kiselev2024} Kiselev N. S., Grabovoy A. V. Unraveling the Hessian: A Key to Smooth Convergence in Loss Function Landscapes // Doklady Mathematics. 2024.
\bibitem{Meshkov2024} Meshkov V., Kiselev N., Grabovoy A. ConvNets Landscape Convergence: Hessian-Based Analysis of Matricized Networks // ISPRAS Open Conference. 2024.

\end{thebibliography}

\end{document}
